\documentclass[12pt,a4paper]{article}

% =====================================================
% PACOTES
% =====================================================
\usepackage[utf8]{inputenc}
\usepackage[T1]{fontenc}
\usepackage[brazil]{babel}
\shorthandoff{"}
\usepackage{geometry}
\usepackage{graphicx}
\usepackage{fancyhdr}
\usepackage{titlesec}
\usepackage[table]{xcolor}
\usepackage{hyperref}
\usepackage{setspace}
\usepackage{enumitem}
\usepackage{newtxtext,newtxmath}
\usepackage[most]{tcolorbox}
\usepackage{float}

\geometry{
    top=2.5cm,
    bottom=3cm,
    left=1.5cm,
    right=1.5cm
}

\onehalfspacing

% =====================================================
% CORES PERSONALIZADAS
% =====================================================
\definecolor{azulTitulo}{RGB}{30,60,90}
\definecolor{azulObjetivos}{RGB}{30,60,90}
\definecolor{cinzaCodigo}{RGB}{245,245,245}
\definecolor{vermelhoCodigo}{RGB}{180,50,50}
\definecolor{mostarda}{RGB}{204,153,0}

% =====================================================
% BOXES
% =====================================================
\newtcolorbox{boxRevisao}[1][]{
    breakable,
    enhanced,
    colback=azulTitulo!5,
    colframe=azulTitulo,
    boxrule=0.5pt,
    arc=4pt,
    left=10pt,
    right=10pt,
    top=10pt,
    bottom=10pt,
    title=#1,
    coltitle=white,
    colbacktitle=azulTitulo,
    fonttitle=\bfseries
}

\newtcolorbox{boxDestaque}{
    colback=mostarda!10,
    colframe=mostarda,
    boxrule=1pt,
    arc=4pt,
    left=8pt,
    right=8pt,
    top=8pt,
    bottom=8pt
}

% =====================================================
% CABEÇALHO E RODAPÉ
% =====================================================
\pagestyle{fancy}
\fancyhf{}
\fancyhead[L]{
    \IfFileExists{../uc-topicos-avancados-latex/figures/garopaba_horizontal_marca2015_PNG.png}{
        \includegraphics[height=1.4cm]{../uc-topicos-avancados-latex/figures/garopaba_horizontal_marca2015_PNG.png}
    }{IFSC - Garopaba}
}
\fancyhead[R]{
    \textbf{Guia de Revisão - AV1}\\
    UC – Tópicos Avançados em Informática\\
    Técnico Integrado em Informática
}
\fancyfoot[L]{\small\textit{Preparação para Avaliação Teórica 01}}
\fancyfoot[R]{\small Página \thepage}
\renewcommand{\footrulewidth}{2.4pt}
\setlength{\headheight}{45pt}

% =====================================================
% ESTILO DE SEÇÃO
% =====================================================
\titleformat{\section}
{\normalfont\Large\bfseries\color{azulTitulo}}
{\thesection}
{1em}
{}
[\vspace{0.3cm}\hrule]

\begin{document}

\begin{center}
    {\LARGE \textbf{Material de Revisão para Avaliação Teórica 01}}\\
    \vspace{0.2cm}
    \textit{Conteúdo: Roteiros 01 ao 07}
\end{center}

\vspace{0.5cm}

\section{Versionamento de Código (Git \& GitHub)}

O Git é essencial para o trabalho colaborativo e segurança do código.

\begin{itemize}
    \item \textbf{Commit:} Salva o progresso localmente com uma mensagem.
    \item \textbf{Push/Pull:} Sincroniza o trabalho com o servidor remoto (GitHub).
    \item \textbf{Clone:} Baixa um projeto completo pela primeira vez.
\end{itemize}

\section{Orientação a Objetos (Python)}

A OO organiza o código em "objetos" que possuem dados e comportamentos.

\begin{boxRevisao}[Conceitos Chave OO]
\begin{itemize}
    \item \textbf{Classe:} O projeto/molde (ex: \texttt{Class Pessoa}).
    \item \textbf{Objeto:} A instância criada (ex: \texttt{joao = Pessoa()}).
    \item \textbf{Atributos:} Variáveis dentro da classe (ex: \texttt{self.nome}).
    \item \textbf{Métodos:} Funções dentro da classe (ex: \texttt{def falar(self)}).
    \item \textbf{Construtor (\texttt{\_\_init\_\_}):} Define o estado inicial do objeto.
\end{itemize}
\end{boxRevisao}

\section{Desenvolvimento Web com Flask}

O Flask é um micro-framework que facilita a criação de servidores web.

\begin{itemize}
    \item \textbf{Rotas:} Endereços que o servidor reconhece (\texttt{@app.route('/')}).
    \item \textbf{Templates:} Arquivos HTML que recebem dados do Python via Jinja2.
    \item \textbf{GET vs POST:} GET para leitura (dados na URL); POST para envio/escrita (dados ocultos no corpo).
    \item \textbf{render\_template:} Função que envia o HTML para o navegador.
\end{itemize}

\section{Virtualização e Docker}

A grande mudança de paradigma no Roteiro 06.

\begin{boxDestaque}
\textbf{VM vs Container:} Enquanto a VM simula o hardware e um SO completo (lento), o Container isola apenas o necessário para a aplicação rodar, compartilhando o "motor" (kernel) do hospedeiro (rápido).
\end{boxDestaque}

\begin{itemize}
    \item \textbf{Imagem:} O arquivo de instalação estático.
    \item \textbf{Container:} O processo em execução.
    \item \textbf{Port Mapping (\texttt{-p 80:80}):} Abre uma janela para o mundo externo acessar o container.
    \item \textbf{Volumes (\texttt{-v}):} Permite editar o código no computador e ver a mudança no container sem reiniciá-lo.
\end{itemize}

\section{Dockerfile e Orquestração}

\begin{boxRevisao}[Dockerfile: A Receita]
\begin{itemize}
    \item \texttt{FROM}: Qual sistema base?
    \item \texttt{WORKDIR}: Onde o código vai morar?
    \item \texttt{COPY}: Quais arquivos levar?
    \item \texttt{RUN}: O que instalar?
    \item \texttt{CMD}: Como ligar o motor?
\end{itemize}
\end{boxRevisao}

\textbf{Docker Compose:} Permite gerenciar um "prédio" de containers (App + Banco) com um único arquivo \texttt{.yml} e o comando \texttt{docker-compose up}.

\end{document}
